\section{Introduction}\label{sec:introduction}

Universal optimality asks whether one configuration minimizes an entire
class of pair interactions.  In the plane, the triangular lattice is the
natural candidate.  The problem goes beyond comparison with other lattices:
a competing configuration may be nonperiodic, may contain arbitrarily close
pairs, and may have very uneven local point counts.  We prove the energy
comparison for this full class under a single centered-density hypothesis.

Put $b=\sqrt3/2$ and let
\[
 A=b^{-1/2}\{j(1,0)+k(1/2,b):j,k\in\Z\}
\]
be the triangular lattice of covolume one.  Let $B_R$ be the closed disk of
radius $R$ centered at the origin.  For a locally finite configuration
$\mathcal C\subset\R^2$, write
$\mathcal C_R=\mathcal C\cap B_R$ and $N_R=\#\mathcal C_R$.  Its centered
disk density is one if
\[
 \frac{N_R}{\pi R^2}\longrightarrow1\qquad(R\to\infty).
\]
For a nonnegative interaction $g:(0,\infty)\to[0,\infty)$, define
\begin{equation}\label{eq:energy-definition}
 E_g(\mathcal C)=\liminf_{R\to\infty}\frac1{N_R}
 \sum_{\substack{x,y\in\mathcal C_R\\x\ne y}}g(\abs{x-y}^2).
\end{equation}
The sum counts ordered pairs.  Density one makes $N_R$ positive for all
sufficiently large $R$, so the lower limit is well defined in $[0,\infty]$.

\Needspace{8\baselineskip}
\begin{theorem}[Universal energy minimality]\label{thm:universal}
Let $g:(0,\infty)\to[0,\infty)$ be smooth and completely monotone, meaning
that $(-1)^r g^{(r)}(t)\ge0$ for every integer $r\ge0$ and every $t>0$.
For every locally finite configuration $\mathcal C\subset\R^2$ of centered
disk density one, the following equality and inequality hold in the extended
nonnegative reals:
\begin{equation}\label{eq:main-inequality}
 E_g(\mathcal C)\ge
 \sum_{a\in A\setminus\{0\}}g(\abs a^2)=E_g(A).
\end{equation}
\end{theorem}

This includes every Gaussian $g(t)=e^{-\pi\alpha t}$ with $\alpha>0$ and
every inverse power $g(t)=t^{-p}$ with $p>0$.  The theorem permits singular
potentials and divergent lattice sums, and it imposes no separation or local
occupancy bound on $\mathcal C$.  It identifies the minimum value, not all
minimizers.  The comparison uses the centered-density and lower-energy
formulation of Cohn et~al.~\citep[Definitions~1.1--1.3]{CKMRV2022}.

\subsection{The planar problem and its predecessors}

The lattice problem began with inverse-power interactions.  For the
convergent Epstein zeta sums, Rankin~\citep{Rankin1953} proved the triangular
minimum in a restricted exponent range, and Cassels~\citep{Cassels1959} addressed
the remaining range.  Cassels subsequently acknowledged a gap in his proof
\citep{Cassels1963Corrigendum}; Ennola~\citep{Ennola1964Lemma} supplied a repair,
and Diananda~\citep{Diananda1964} simplified the needed lemmas.  Together these works
establish the triangular minimum among planar lattices of fixed covolume
for $g(t)=t^{-p}$ with $p>1$.  Montgomery's theta theorem gives the
Gaussian minimum for every positive parameter
\citep[Theorem~1]{Montgomery1988}.  These results identify the best lattice;
they leave comparison with nonlattice configurations open.

Cohn and Kumar placed this question alongside their study of universal
energy minimization on spheres
\citep[Theorem~1.2]{CohnKumar2007}.  Their Euclidean conjecture singled out
the triangular, $E_8$, and Leech lattices
\citep[Conjecture~9.4]{CohnKumar2007}.  It asks for a sharp auxiliary function
for each sufficiently decaying completely monotone squared-distance
interaction: a Fourier lower bound that attains the lattice energy and
therefore proves the comparison for periodic competitors.  In dimensions
$8$ and $24$, Cohn, Kumar, Miller, Radchenko, and Viazovska proved universal
optimality and constructed sharp Gaussian bounds and Fourier interpolation
formulas \citep[Theorems~1.4, 1.7, and 1.9]{CKMRV2022}.

Several planar results address restricted competitor classes.
Faulhuber et~al.~\citep[Theorem~1.1 and Corollary~1.3]{FaulhuberShafkulovskaZlotnikov2023}
compare the triangular lattice with the honeycomb configuration,
Cartesian products of one-dimensional periodic configurations, and unions
of two translates of a rectangular or triangular lattice, with each class
normalized to the same density.
Hardin and Tenpas~\citep[Theorem~1]{HardinTenpas2025} prove triangular optimality among
four-point configurations modulo a fixed triangular lattice and six-point
configurations modulo $\Z\times\sqrt3\Z$, with the appropriate scaling
and rotation of the minimizing lattice.  Their completely monotone
interactions satisfy a rapid-decay assumption.
Lebl\'e~\citep[Theorems~1--2]{Leble2025} proves that sufficiently small bounded
displacements of the triangular lattice cannot lower the energy; the
permitted displacement size depends on the interaction.
These restrictions are different from the centered-density hypothesis of
Theorem~\ref{thm:universal}.

Theorem~\ref{thm:universal} proves the positive planar universal-energy
conclusion in the formulation stated above.  The construction also supplies
sharp auxiliary functions for every Gaussian.  Integrating the resulting
energy inequalities does not construct a sharp auxiliary for each mixed
potential in the class required by Cohn and Kumar's conjecture.  That
stronger per-potential quantifier is not claimed here.

\subsection{The Gaussian route and the interpolation problem}

The proof reduces the energy problem to a family of Fourier inequalities.
We use the normalization
\[
 \widehat f(\xi)=\int_{\R^2}f(x)e^{-2\pi i x\cdot\xi}\,\dd x
\]
and write
$A^*=\{w\in\R^2:w\cdot a\in\Z\text{ for every }a\in A\}$
for the dual lattice.  Fix a Gaussian
$G_\alpha(x)=e^{-\pi\alpha|x|^2}$, with $\alpha>0$.
We seek a real radial Schwartz function $f_\alpha$ satisfying
\[
 f_\alpha\le G_\alpha,\qquad \widehat f_\alpha\ge0,
 \qquad
 \begin{aligned}
 f_\alpha(a)&=G_\alpha(a)&& (a\in A\setminus\{0\}),\\
 \widehat f_\alpha(w)&=0&& (w\in A^*\setminus\{0\}).
 \end{aligned}
\]
The two inequalities give a lower bound for arbitrary competitors.  The
contacts at the direct and dual lattice points make the bound exact on $A$.
For every competitor $\mathcal C$ of centered disk density one,
Section~\ref{sec:energy} proves
\[
 E_{e^{-\pi\alpha t}}(\mathcal C)
 \ge\widehat f_\alpha(0)-f_\alpha(0)
 =\sum_{a\in A\setminus\{0\}}e^{-\pi\alpha|a|^2}.
\]
The lower bound compares the point measure in a disk with area measure in
a slightly larger disk.  Its error is uniform even when points cluster.
The equality follows from Poisson summation on $A$; the competitor itself
need not be periodic.

Exact contacts and signs between contacts are different requirements.  The
first is an infinite interpolation problem.  The second must hold at every
radius and every Gaussian scale.  Our construction separates them, then
uses a quantitative correction estimate to bring them together.
The normalized auxiliary theorem, Theorem~\ref{thm:normalized-pair}, is
stated after the shell and node definitions in Section~\ref{sec:fourier}.
Its periodic set of rescaled squared radii contains all triangular-lattice
shells, so the prescribed values and first derivatives (jets) supply the
required contacts.

The construction has three stages.  First, Section~\ref{sec:fourier}
encodes the periodic contact set by a sine product with double zeros.
Multiplying that product by summable double-pole and simple-pole terms gives
entire functions with independently specified values and slopes at the
nodes.  A compactly supported spectral measure represents each such cardinal
function.  Gaussian damping turns it into a radial Schwartz function, and transforming
the Gaussian inside the integral gives its genuine Fourier transform.
Combining two cardinal functions with their transformed profiles reduces
the simultaneous contacts to coupled linear equations for two summable
coefficient lists.  The transformed jets become uniformly small at distant
output nodes, regardless of the input coefficient locations.

Second, Section~\ref{sec:finite} constructs finite reference lists by
positive quadrature, leaving the spectral atoms exact.  It proves a uniform
quadrature error and certifies explicit matrix and polynomial inequalities.
These lists do not yet interpolate exactly.  The available bound on the
full coupled operator is also too large for a direct small-perturbation
argument.  Section~\ref{sec:interpolation} instead inverts the finite block
and eliminates it; the small distant outputs make the remaining tail
operator invertible.  This solves all the exact interpolation equations and
bounds the correction to the reference lists.

Third, Section~\ref{sec:signs} proves the signs corresponding to
$G_\alpha-f_\alpha\ge0$ and $\widehat f_\alpha\ge0$.  It works in the
rescaled squared-radius coordinate, after removing positive Gaussian factors.
The normalized difference and Fourier profile then vanish to second order
at the positive nodes because of the exact value-and-slope contacts.  To control them
near a node, we subtract the value and linear term of each component, divide
by the squared distance to the node, and then approximate the quotient by a
polynomial.  This keeps a small interpolation error from being divided by a vanishing square.
Coefficients in the Bernstein basis bound the polynomial on whole intervals
between nodes, and rectangular bounds on its parameter data cover whole
parameter intervals.  On the unbounded radial tail, the signed contribution
from the low-node rational terms is positive.  The sum of the remaining
terms has double zeros and a small second derivative.  A quadratic lower
bound for the sine product puts the positive part and the remainder on the
same scale.  Appendix~\ref{sec:arithmetic} supplies the scalar bounds used
in this argument and describes the accompanying interval verifier for the
exact quadrature arrays.  Appendix~\ref{app:floor-certificate} gives a
separate rational verification procedure, whose implication is conditional
on its tests passing.

Section~\ref{sec:energy} converts the normalized construction into the
functions $f_\alpha$ above.  It directly covers $\alpha\ge1$; a
Fourier-complement construction supplies $0<\alpha<1$.
Finally, Section~\ref{sec:mixtures} represents each positive shift of a
completely monotone potential as
\[
 g(\varepsilon+t)=\int_{[0,1]}v^t\,\dd\rho_\varepsilon(v)
 \qquad(t>0),
\]
where $\varepsilon>0$ and $\rho_\varepsilon$ is positive with total mass
$g(\varepsilon)$.  The shift makes this mass finite even when $g$ is
singular at zero.  For $0<v<1$, the kernel $v^t$ is a Gaussian with
$\alpha=-\log(v)/\pi$; the endpoint $v=1$ gives the constant interaction.
Fatou's lemma and Tonelli's theorem transfer the Gaussian comparisons to
the mixture, and the shift is removed only in the lattice sum.  This proves
Theorem~\ref{thm:universal}, including singular potentials and infinite
energies.

\subsection{Antecedents of the method}

The interpolation mechanism has substantial antecedents in the
Beurling--Selberg extremal method.  Cardinal formulas with squared
trigonometric factors, double-pole value terms, and simple-pole derivative
terms occur in the one-dimensional construction of Cohn and Kumar~\citep[Proposition~9.6]{CohnKumar2007} and in the Gaussian
extremal functions of Carneiro et~al.~\citep[Section~1, (1.2)--(1.3)]{CarneiroLittmannVaaler2013}.
Radchenko and Viazovska~\citep[Theorem~1]{RadchenkoViazovska2019} reconstruct even Schwartz functions
on the real line from simultaneous function and Fourier data at square-root
nodes.  Theorems~1.7 and~1.9 of Cohn et~al.~\citep{CKMRV2022} give radial
value-and-derivative reconstruction formulas in dimensions $8$ and $24$.
In the plane, simultaneous values and first derivatives of a radial
Schwartz function and its Fourier transform on the triangular shells do
not determine the function: Talebizadeh Sardari~\citep[Theorem~1.7]{Sardari2021} proves that
infinitely many linearly independent functions have all these data zero.
His proof uses periodic congruence supersets of the triangular
quadratic-form values \citep[Section~6.1, Lemma~6.1]{Sardari2021}.
Our construction prescribes data on a periodic superset and proves
uniqueness within its specified summable coefficient system.  The following
sections establish exact interpolation in that system and the
parameter-uniform planar signs.

Related Fourier sign conditions enter the sphere-packing bound of
Cohn and Elkies~\citep{CohnElkies2003}.  The companion manuscript
\citep[Sections~2--4]{SharpPlanar2026} develops the cardinal,
Gaussian-transform, and finite sign-certification framework for a sharp planar
packing bound.  We adapt that framework to a continuously varying Gaussian
target and establish the required energy identities, correction estimates,
and signs locally.  The separate atomic companion independently proves the same
energy conclusion \mbox{\citep[Theorem~1.1]{AtomicTriangular2026}}; this is not a
premise of the present proof.  Its auxiliary construction is compared with
ours after Theorem~\ref{thm:normalized-pair}.

The density-only linear-programming argument follows
Cohn and de~Courcy-Ireland~\citep[\mbox{Proposition~2.2}]{CohnDeCourcyIreland2018}, whose background-measure
method follows Cohn and Zhao~\citep[Theorem~3.3]{CohnZhao2014}.  The latter is a packing
theorem; we use the method, not its exclusion-distance hypotheses.
Section~\ref{sec:energy} proves the needed statement for real even Schwartz
functions with nonnegative Fourier transform.  The positive-mixture
representation belongs to the classical Hausdorff--Bernstein--Widder theory
\citep{Hausdorff1921,Bernstein1929,Widder1931}; our finite-difference and
moment-measure proof is given in Section~\ref{sec:mixtures}.  The lower-limit,
endpoint, and infinite-energy arguments are proved there as well.

\subsection{Renormalized and jellium consequences}

For $0<s<2$, the raw Riesz energy with $g(t)=t^{-s/2}$ is infinite even
on $A$, so the comparison of raw energies does not by itself distinguish
configurations in this range.  Neutralized energies account for a uniform
oppositely charged background.  In a finite jellium system, one combines
point--point, point--background, and background--background interactions;
the thermodynamic jellium minimum is the large-volume limit of the minimum
energy per unit volume.  Renormalized energies remove the divergent
self-energy of each point.  Their periodic version compares configurations
in a fixed torus, and their infinite-system version measures field energy
per unit area.  Petrache and Serfaty~\citep[Theorem~2]{PetracheSerfaty2020} proved that the
Cohn--Kumar universal-optimality conjecture implies the corresponding
Coulomb and Riesz crystallization results.  Their heat-kernel representation
\citep[Lemma~1]{PetracheSerfaty2020} supplies the connection to Gaussian
comparisons used here.  We carry out this passage for the field functional
defined in Section~\ref{sec:renormalized}, including its normalization and
periodic approximation, after completing the main theorem in
Section~\ref{sec:mixtures}.

Corollary~\ref{cor:riesz-jellium} proves the triangular minimum for every
$0\le s<2$, with interaction $-\log|x|$ at $s=0$ and $|x|^{-s}$ for
$0<s<2$. On each torus $\R^2/(nA)$, the $n^2$ classes of $A/(nA)$
minimize the canonical periodic energy among all configurations of $n^2$
distinct points. The lattice also minimizes the infinite-system field
energy $W_s$ among density-one configurations, and the ordered-pair
thermodynamic jellium minimum is
$W_s(A)/\kappa_s$, where $\kappa_0=2\pi$ and $\kappa_s=4\pi$ for
$0<s<2$.

Section~\ref{sec:renormalized} defines these quantities with unit background
and fixes the field normalization by the fundamental-solution equation.
The field energy takes the centered-square upper limit at each fixed
self-energy cutoff, then removes the cutoff; the infimum over compatible
gradient fields is outside both limits. The proof compares heat-kernel
energies for arbitrary period lattices. Appendix~\ref{app:normalized-periodic-approximation}
constructs a sequence of square-periodic fields approaching the field
infimum, so the comparison does not require a periodic approximation with
prescribed triangular periods. The scalar periodic and jellium identities
of Lewin et~al.~\citep[Section~VI, Theorem~2]{LewinLiebSeiringer2019} and
Lauritsen~\citep[Theorem~II.1]{Lauritsen2021} then identify the thermodynamic value.

The fresh long-range conclusion here is $0<s<2$.  At $s=0$, the corollary
overlaps the independent direct logarithmic result of the Coulomb companion
\citep[Theorems~1.1--1.2]{CoulombRenormalized2026}, which uses its own field-energy
normalization and proves a lower bound on square tori.  The periodic assertion
above concerns only the compatible triangular tori $\R^2/(nA)$.  It asserts
neither uniqueness nor exclusion of defects, and makes no claim at $s=2$ or
for quantum crystallization.
